\documentclass[12pt]{article}
\usepackage{amsmath, amssymb, geometry, enumitem}
\geometry{letterpaper, margin=1in}

\title{Math Practice Set: Equations of Lines and Planes}
\author{Based on Homework 12.5 Concepts}
\date{}

\begin{document}

\maketitle

\section*{Concept Checklist}
This problem set tests the following concepts found in the reference material:
\begin{itemize}
    \item \textbf{Lines:} Vector/Parametric/Symmetric forms, Segments, Zero-components.
    \item \textbf{Planes:} Normal vectors, 3-Point derivation, Intercepts.
    \item \textbf{Relationships:} Parallelism, Perpendicularity, Skew lines, Angles.
    \item \textbf{Intersections:} Line-Line, Line-Plane, Plane-Plane.
    \item \textbf{Applications:} Distance formulas, Finance (Portfolio) models, Error analysis.
\end{itemize}

\section*{Part 1: Practice Problems}

\subsection*{Conceptual Questions}
\textbf{1.} Determine if the statement is True or False: Two lines in $\mathbb{R}^3$ that are not parallel must intersect.

\textbf{2.} Determine if the statement is True or False: If a line is perpendicular to a plane, it is parallel to the plane's normal vector.

\subsection*{Line Equations}
\textbf{3.} Find the vector equation and parametric equations for the line passing through $P(2, -1, 5)$ and parallel to $\mathbf{v} = \langle 4, 3, -7 \rangle$.

\textbf{4.} Find the symmetric equations for the line passing through $A(1, 3, 2)$ and $B(-4, 3, 5)$.

\textbf{5.} Find parametric equations for the line passing through $(6, 2, 1)$ that is perpendicular to the plane $4x - 2y + z = 10$.

\textbf{6.} Find the vector equation for the line segment connecting $P(1, 0, 1)$ to $Q(3, -4, 2)$. Specify the range for the parameter $t$.

\textbf{7.} A line passes through the origin and the point $(5, -2, 8)$. Write its symmetric equations.

\textbf{8.} Write the symmetric equations for the line through $(4, -1, 3)$ parallel to $\mathbf{v} = \langle 2, 0, -5 \rangle$.

\subsection*{Line Relationships}
\textbf{9.} Determine if the lines $L_1$ and $L_2$ are parallel, perpendicular, or neither. \\
$L_1: x = 2 - 3t, y = 4 + t, z = 6t$ \\
$L_2: \frac{x+1}{6} = \frac{y-2}{-2} = \frac{z}{-12}$

\textbf{10.} Determine if the lines $L_1$ and $L_2$ are parallel, perpendicular, or neither. \\
$L_1: \mathbf{r}(t) = \langle 1, 0, 2 \rangle + t\langle 3, -1, 2 \rangle$ \\
$L_2: \mathbf{r}(s) = \langle 4, 1, 5 \rangle + s\langle 1, 5, 1 \rangle$

\textbf{11.} Determine if the lines intersect or are skew. If they intersect, find the point. \\
$L_1: x = 1 + t, y = -2 + 3t, z = 4 - t$ \\
$L_2: x = 2s, y = 3 + s, z = -3 + 4s$

\subsection*{Plane Equations}
\textbf{12.} Find the parametric equations for the line of intersection of the planes $x + y + z = 1$ and $x - y + z = 0$.

\textbf{13.} Find the scalar equation of the plane passing through $(2, 4, -1)$ with normal vector $\mathbf{n} = \langle 3, 1, 5 \rangle$.

\textbf{14.} Find the equation of the plane passing through $(-1, 2, 3)$ that is parallel to the plane $2x - 3y + z = 6$.

\textbf{15.} Find the linear equation ($ax+by+cz=d$) of the plane passing through the points $(1, 0, 0)$, $(0, 2, 0)$, and $(0, 0, 3)$.

\textbf{16.} Find the equation of the plane that contains the point $(1, -1, 2)$ and the line $x = 2t, y = 3 + t, z = 4 - 2t$.

\textbf{17.} Find the equation of the plane passing through $(4, 1, 1)$ that is perpendicular to both planes $x + y + z = 1$ and $2x + y - 3z = 0$.

\textbf{18.} Find the x, y, and z intercepts of the plane $3x - 4y + 2z = 24$.

\subsection*{Intersections and Angles}
\textbf{19.} Find the point where the line $x = 2 - t, y = 1 + 3t, z = 4t$ intersects the plane $2x + y - z = 2$.

\textbf{20.} Determine if the line $\mathbf{r}(t) = \langle 1, 2, 3 \rangle + t\langle 2, -1, 3 \rangle$ is parallel to, contained in, or intersects the plane $x + 5y + z = 7$.

\textbf{21.} Find the angle (in degrees, rounded to 1 decimal) between the planes $x + y = 1$ and $2x + y - 2z = 2$.

\textbf{22.} Find the distance from the point $(1, 1, 1)$ to the plane $x + 2y + 2z = 0$.

\subsection*{Advanced Concept Checks}
\textbf{23.} "Find the Flaw": A student claims the direction vector of the line given by $\frac{x-1}{2} = \frac{2-y}{3} = z$ is $\langle 2, 3, 1 \rangle$. Is this correct? If not, find the correct direction.

\textbf{24.} Find the equation of the plane that contains the two intersecting lines: \\
$L_1: x = 1 + t, y = 1 - t, z = 2t$ \\
$L_2: x = 1 - s, y = 1 + 2s, z = 3s$

\textbf{25.} Find the parametric equations for the line through $(1, 0, 1)$ that is parallel to the planes $x + y + z = 5$ and $3x - y = 4$.

\textbf{26.} Find the equation of the plane that contains the line $x = 1 + t, y = 2 - t, z = 3t$ and is perpendicular to the plane $3x - 2y + z = 4$.

\subsection*{Applications}
\textbf{27. Finance (Portfolio Optimization):} 
Three assets have weights $w_1, w_2, w_3$. 
The budget constraint is a plane: $w_1 + w_2 + w_3 = 1$.
The return constraint is a plane: $0.05w_1 + 0.10w_2 + 0.15w_3 = 0.10$ (targeting 10\% return).
The valid portfolios lie on the line of intersection of these two planes. Find the vector equation of this line.

\textbf{28.} Find the distance between the parallel planes $2x - 2y + z = 3$ and $2x - 2y + z = 9$.

\textbf{29.} Find the intersection of the line passing through $(1,0,1)$ and $(2,1,0)$ with the $xz$-plane.

\textbf{30.} Comprehensive Synthesis: 
Consider the points $A(1,0,0)$, $B(0,1,0)$, and $C(0,0,1)$.
(a) Find the equation of the plane containing $A, B, C$.
(b) Find the line passing through the origin perpendicular to this plane.
(c) Find the intersection point of this line and the plane.

\newpage
\section*{Part 2: Detailed Solutions}

\textbf{1. False.} Lines in $\mathbb{R}^3$ can be skew (neither parallel nor intersecting).

\textbf{2. True.} The definition of a line perpendicular to a plane implies its direction vector is parallel to the plane's normal.

\textbf{3.} 
Point $P_0(2, -1, 5)$, $\mathbf{v} = \langle 4, 3, -7 \rangle$. \\
Vector: $\mathbf{r} = \langle 2, -1, 5 \rangle + t\langle 4, 3, -7 \rangle$. \\
Parametric: $x = 2 + 4t, y = -1 + 3t, z = 5 - 7t$.

\textbf{4.} 
$\mathbf{v} = B - A = \langle -4-1, 3-3, 5-2 \rangle = \langle -5, 0, 3 \rangle$. \\
Since the y-component is 0, we separate it. \\
Symmetric: $\frac{x-1}{-5} = \frac{z-2}{3}, \quad y = 3$.

\textbf{5.} 
Plane normal $\mathbf{n} = \langle 4, -2, 1 \rangle$. The line is perpendicular to the plane, so $\mathbf{v} = \mathbf{n} = \langle 4, -2, 1 \rangle$. \\
Parametric: $x = 6 + 4t, y = 2 - 2t, z = 1 + t$.

\textbf{6.} 
$\mathbf{r}_0 = \langle 1, 0, 1 \rangle$. $\mathbf{v} = Q - P = \langle 2, -4, 1 \rangle$. \\
$\mathbf{r}(t) = \langle 1, 0, 1 \rangle + t\langle 2, -4, 1 \rangle$ for $0 \le t \le 1$.

\textbf{7.} 
Origin $O(0,0,0)$, $P(5,-2,8)$. $\mathbf{v} = \langle 5, -2, 8 \rangle$. \\
Symmetric: $\frac{x}{5} = \frac{y}{-2} = \frac{z}{8}$.

\textbf{8.} 
Point $(4, -1, 3)$, $\mathbf{v} = \langle 2, 0, -5 \rangle$. \\
Symmetric: $\frac{x-4}{2} = \frac{z-3}{-5}, \quad y = -1$.

\textbf{9.} 
$\mathbf{v}_1 = \langle -3, 1, 6 \rangle$. $\mathbf{v}_2 = \langle 6, -2, -12 \rangle$. \\
Ratio check: $\frac{6}{-3} = -2$, $\frac{-2}{1} = -2$, $\frac{-12}{6} = -2$. \\
Since ratios are constant, $\mathbf{v}_2 = -2\mathbf{v}_1$. They are \textbf{Parallel}.

\textbf{10.} 
$\mathbf{v}_1 = \langle 3, -1, 2 \rangle$. $\mathbf{v}_2 = \langle 1, 5, 1 \rangle$. \\
Dot product: $\mathbf{v}_1 \cdot \mathbf{v}_2 = (3)(1) + (-1)(5) + (2)(1) = 3 - 5 + 2 = 0$. \\
Since the dot product is zero, they are \textbf{Perpendicular}.

\textbf{11.} 
$x: 1+t = 2s \implies t = 2s - 1$. \\
$y: -2+3t = 3+s$. Substitute $t$: $-2 + 3(2s-1) = 3+s \implies -2 + 6s - 3 = 3+s \implies 5s = 8 \implies s = 1.6$. \\
$t = 2(1.6) - 1 = 2.2$. \\
Check $z$: $L_1(z) = 4 - 2.2 = 1.8$. $L_2(z) = -3 + 4(1.6) = -3 + 6.4 = 3.4$. \\
$1.8 \neq 3.4$. They do not intersect. Since directions are not parallel, they are \textbf{Skew}.

\textbf{12.} 
Normals $\mathbf{n}_1 = \langle 1, 1, 1 \rangle, \mathbf{n}_2 = \langle 1, -1, 1 \rangle$. \\
Direction $\mathbf{v} = \mathbf{n}_1 \times \mathbf{n}_2 = \begin{vmatrix} \mathbf{i} & \mathbf{j} & \mathbf{k} \\ 1 & 1 & 1 \\ 1 & -1 & 1 \end{vmatrix} = \langle (1-(-1)), -(1-1), (-1-1) \rangle = \langle 2, 0, -2 \rangle$. \\
Find point: Set $z=0$. $x+y=1, x-y=0 \implies 2x=1 \implies x=1/2, y=1/2$. Point $(1/2, 1/2, 0)$. \\
Parametric: $x = 0.5 + 2t, y = 0.5, z = -2t$. (Simplified direction $\langle 1, 0, -1 \rangle$ is also valid).

\textbf{13.} 
$3(x-2) + 1(y-4) + 5(z - (-1)) = 0$. \\
$3x - 6 + y - 4 + 5z + 5 = 0 \implies 3x + y + 5z = 5$.

\textbf{14.} 
Normal of given plane $\mathbf{n} = \langle 2, -3, 1 \rangle$. Parallel plane has same normal. \\
$2(x+1) - 3(y-2) + 1(z-3) = 0$. \\
$2x + 2 - 3y + 6 + z - 3 = 0 \implies 2x - 3y + z = -5$.

\textbf{15.} 
Vectors in plane: $\mathbf{a} = P_2 - P_1 = \langle -1, 2, 0 \rangle$, $\mathbf{b} = P_3 - P_1 = \langle -1, 0, 3 \rangle$. \\
$\mathbf{n} = \mathbf{a} \times \mathbf{b} = \langle 6, -(-3), 2 \rangle = \langle 6, 3, 2 \rangle$. \\
$6(x-1) + 3(y-0) + 2(z-0) = 0 \implies 6x + 3y + 2z = 6$.

\textbf{16.} 
Plane contains point $P(1, -1, 2)$ and line. Line has point $L_0(0, 3, 4)$ and direction $\mathbf{v} = \langle 2, 1, -2 \rangle$. \\
Vector in plane $\mathbf{u} = P - L_0 = \langle 1, -4, -2 \rangle$. \\
$\mathbf{n} = \mathbf{v} \times \mathbf{u} = \langle 2, 1, -2 \rangle \times \langle 1, -4, -2 \rangle$. \\
$\mathbf{n} = \langle -2-8, -(-4+2), -8-1 \rangle = \langle -10, 2, -9 \rangle$. \\
Equation: $-10(x-1) + 2(y+1) - 9(z-2) = 0 \implies -10x + 10 + 2y + 2 - 9z + 18 = 0$. \\
$-10x + 2y - 9z = -30 \implies 10x - 2y + 9z = 30$.

\textbf{17.} 
$\mathbf{n} = \mathbf{n}_1 \times \mathbf{n}_2$ where $\mathbf{n}_1 = \langle 1, 1, 1 \rangle, \mathbf{n}_2 = \langle 2, 1, -3 \rangle$. \\
$\mathbf{n} = \langle -3-1, -(-3-2), 1-2 \rangle = \langle -4, 5, -1 \rangle$. \\
Equation: $-4(x-4) + 5(y-1) - 1(z-1) = 0 \implies -4x + 16 + 5y - 5 - z + 1 = 0$. \\
$-4x + 5y - z = -12 \implies 4x - 5y + z = 12$.

\textbf{18.} 
x-int ($y=z=0$): $3x=24 \implies x=8$. \\
y-int ($x=z=0$): $-4y=24 \implies y=-6$. \\
z-int ($x=y=0$): $2z=24 \implies z=12$.

\textbf{19.} 
Sub line into plane: $2(2-t) + (1+3t) - (4t) = 2$. \\
$4 - 2t + 1 + 3t - 4t = 2 \implies -3t + 5 = 2 \implies -3t = -3 \implies t = 1$. \\
Point: $x=2-1=1, y=1+3=4, z=4$. Point $(1, 4, 4)$.

\textbf{20.} 
Check orthogonality of $\mathbf{v}$ and $\mathbf{n}$. $\mathbf{v} = \langle 2, -1, 3 \rangle, \mathbf{n} = \langle 1, 5, 1 \rangle$. \\
$\mathbf{v} \cdot \mathbf{n} = 2(1) - 1(5) + 3(1) = 2 - 5 + 3 = 0$. \\
Line is parallel or contained. Check point $(1, 2, 3)$ in plane: $1 + 5(2) + 3 = 14 \neq 7$. \\
Answer: \textbf{Parallel} (No intersection).

\textbf{21.} 
$\mathbf{n}_1 = \langle 1, 1, 0 \rangle, \mathbf{n}_2 = \langle 2, 1, -2 \rangle$. \\
$|\mathbf{n}_1| = \sqrt{2}$. $|\mathbf{n}_2| = \sqrt{4+1+4} = 3$. \\
Dot product: $2+1+0 = 3$. \\
$\cos \theta = \frac{3}{\sqrt{2} \cdot 3} = \frac{1}{\sqrt{2}}$. \\
$\theta = 45^\circ$.

\textbf{22.} 
$D = \frac{|1(1) + 2(1) + 2(1) + 0|}{\sqrt{1^2 + 2^2 + 2^2}} = \frac{|5|}{\sqrt{9}} = \frac{5}{3}$.

\textbf{23.} 
\textbf{Flaw:} The term $\frac{2-y}{3}$ is not in standard form $\frac{y-y_0}{b}$. \\
Rewrite: $\frac{y-2}{-3}$. The term $z$ is $\frac{z-0}{1}$. \\
Correct Direction: $\langle 2, -3, 1 \rangle$.

\textbf{24.} 
Identify directions $\mathbf{v}_1 = \langle 1, -1, 2 \rangle, \mathbf{v}_2 = \langle -1, 2, 3 \rangle$. \\
Intersection point (set $t=0, s=0$): $P(1, 1, 0)$. \\
$\mathbf{n} = \mathbf{v}_1 \times \mathbf{v}_2 = \langle -3-4, -(-3+2), 2-1 \rangle = \langle -7, -5, 1 \rangle$. \\
$-7(x-1) - 5(y-1) + 1(z-0) = 0 \implies -7x + 7 - 5y + 5 + z = 0$. \\
$7x + 5y - z = 12$.

\textbf{25.} 
Line is parallel to planes $\implies \mathbf{v} \perp \mathbf{n}_1$ and $\mathbf{v} \perp \mathbf{n}_2$. \\
$\mathbf{n}_1 = \langle 1, 1, 1 \rangle, \mathbf{n}_2 = \langle 3, -1, 0 \rangle$. \\
$\mathbf{v} = \mathbf{n}_1 \times \mathbf{n}_2 = \langle 0-(-1), -(0-3), -1-3 \rangle = \langle 1, 3, -4 \rangle$. \\
Line: $x = 1 + t, y = 3t, z = 1 - 4t$.

\textbf{26.} 
Plane contains line $\implies \mathbf{n} \perp \mathbf{v}_{line}$. $\mathbf{v}_{line} = \langle 1, -1, 3 \rangle$. \\
Plane is perpendicular to Plane 2 $\implies \mathbf{n} \perp \mathbf{n}_{plane2}$. $\mathbf{n}_{plane2} = \langle 3, -2, 1 \rangle$. \\
$\mathbf{n} = \langle 1, -1, 3 \rangle \times \langle 3, -2, 1 \rangle = \langle -1+6, -(1-9), -2+3 \rangle = \langle 5, 8, 1 \rangle$. \\
Point from line: $(1, 2, 0)$. \\
$5(x-1) + 8(y-2) + 1(z-0) = 0 \implies 5x + 8y + z = 21$.

\textbf{27.} 
System: $w_1 + w_2 + w_3 = 1$ and $w_1 + 2w_2 + 3w_3 = 2$ (after simplifying returns). \\
$\mathbf{n}_1 = \langle 1, 1, 1 \rangle, \mathbf{n}_2 = \langle 1, 2, 3 \rangle$. \\
$\mathbf{v} = \mathbf{n}_1 \times \mathbf{n}_2 = \langle 1, -2, 1 \rangle$. \\
Find point (set $w_3=0$): $w_1 + w_2 = 1, w_1 + 2w_2 = 2 \implies w_2 = 1, w_1 = 0$. Point $(0, 1, 0)$. \\
Vector Eq: $\mathbf{w} = \langle 0, 1, 0 \rangle + t\langle 1, -2, 1 \rangle$.

\textbf{28.} 
Pick point on Plane 1: Let $y=z=0 \implies 2x=3 \implies x=1.5$. $P(1.5, 0, 0)$. \\
Dist to Plane 2 ($2x - 2y + z - 9 = 0$): \\
$D = \frac{|2(1.5) - 0 + 0 - 9|}{\sqrt{2^2 + (-2)^2 + 1^2}} = \frac{|3-9|}{\sqrt{9}} = \frac{6}{3} = 2$.

\textbf{29.} 
Find line through $(1,0,1)$ and $(2,1,0)$. $\mathbf{v} = \langle 1, 1, -1 \rangle$. \\
$x = 1+t, y = t, z = 1-t$. \\
XZ-plane means $y=0$. \\
$t=0 \implies x=1, z=1$. Point $(1, 0, 1)$.

\textbf{30.} 
(a) Normal $\mathbf{n} = \langle 1, 1, 1 \rangle$ (Standard intercept plane). $x+y+z=1$. \\
(b) Line $\perp$ Plane through origin. $\mathbf{v} = \langle 1, 1, 1 \rangle$. $\mathbf{r}(t) = \langle t, t, t \rangle$. \\
(c) Intersection: $t + t + t = 1 \implies 3t=1 \implies t=1/3$. Point $(1/3, 1/3, 1/3)$.

\end{document}